\scanpage{111}
Faisceaux $l$-adiques admissibles et motifs virtuels.

\deleted{1. Soit $S$ un schéma, $F$ un \unsure{$Q_l$|$\mathbb Q_\ell$}}

1. On se fixe un nb premier $l$. Les schémas envisagés sont premiers à $l$. (\uline{Le cas le plus intéressant est celui des schémas} loc. de \uline{type fini sur $\mathbb Z$}.) Les faisceaux envisagés sont les $Q_l$-faisceaux \inserted{constructibles}. On note $Q_l(1)$ le faisceau de Tate $T_l(G_m)\otimes_{\mathbb Z_l}Q_l$, et le twist par sa puissance tensorielle $n$ième est noté $F\longmapsto F(n)$.

% slanted
\margin{se borner pour commencer à ce cas !}
% slanted
\margin{On travaille dans la catégorie des \textquotedblleft faisceaux $l$-adiques de nature arithmétique\ldots}
% slanted
\margin{définition concrète : dire qu'on se ramène à une situation de t.f. sur $\operatorname{Sp}\mathbb Z$}

2.1 Un faisceau $F$ sur $S$ est dit \uline{spécial} \inserted{\uline{de poids} $i$} s'il est isomorphe à un \deleted{facteur direct} \inserted{quotient \deleted{\unsure{constant}} tordu} d'un \inserted{sous} faisceau \inserted{constant tordu d'un faisceau} de la forme $R^i f_*(Q_{lX})$, où $f:X\to S$ est projectif et lisse \inserted{de dim relative $\unsure{i|j}$ \unsure{plutôt ?}}. \unsure{(|[}On notera que si $i<0$, cela implique $F=0$ ; si $i=0$, cela \deleted{\unsure{équivaut} \unsure{implique}} équivaut à : $F$ localement constant.$^{(+)}$

% slanted
\margin{\deleted{\unsure{Attention}, la réciproque n'est pas vraie ; même pour un groupe fini seulement, il faut une repr. sur $Q$, et \unsure{pas que} \illegible{2} \uline{\unsure{définissable sur $Q$}} !}}

\uline{Lemme} 2.\unsure{1|2}. \quad a) Stabilité de \textquotedblleft spécial de poids $i$\textquotedblright{} par \inserted{sous-\unsure{facteurs|quotients} et quotients cst tordus, et} sommes directes.

b) Si $F,F'$ sont spéciaux de poids respectifs $i,j$, alors $F\otimes F'$ est spécial de $i+j$.

c) Stabilité par images inverses.

d) Si $F$ spécial de poids $i$, $\check F(-i)$ l'est aussi

e) $Q_l(-i)$ est spécial de poids $\deleted{\unsure{-|a}}\unsure{2i|i}$ pour $i\ge0$

a) se prouve en introduisant un $\coprod_{1\le i\le n}X_i$.

b) se prouve en considérant un $X_1\times_S X_2$, et utilisant Künneth

c) Par th. de changement de base.

\margin{d) utiliser Lefschetz}

d) \uline{Par th. de Lefschetz} en cohomologie $l$-adique.

e) On \deleted{peut} a $Q_l(-i)=R^{2i}f_*(Q_{l\mathbb P^i_S})$.

\margin{\uline{Cor} 2\unsure{?|3} Si $F$ spécial de poids $i$, et $n\ge i$, alors $\check F(-n)$ est spécial de poids $2n-i$}

3.1 Un faisceau $F$ sur $S$ est dit \uline{effectivement} \uline{admissible} si \inserted{pour tout ouvert affine $U$ de $S$} on peut en trouver une \deleted{filtration de $F$\unsure{$|S$|$|U$} dont les quotients} \inserted{filtration \deleted{de \unsure{$F|U$}} de $F|U$ dont les quotients} successifs sont de la forme $i_!(G)$, où $i:Y\to X$ est une \deleted{\unsure{et de la forme} $i_!(\phantom G)$ \unsure{d'une partie de $U$ en partie}} imm. de présentation finie et où $G$ sur \unsure{$S_Y$|$Y$} est \deleted{\unsure{solution} localement isomorphes \unsure{constructibles} \unsure{pour les $F|S$ sont|pour les $F|S_\alpha$ sont}} admissible.

\margin{il faut se permettre une extension finie surj. loc. libres radicielles de $Y$, pour rendre $G$ spécial}

\deleted{On dit} Si $P$ est une partie de $\mathbb Z$, on dit que $F$ est à poids $\subset P$ si \deleted{les \unsure{$F|S_i$|$F|S_\alpha$} sont} \inserted{on peut prendre les $G$} admissibles de poids $\in P$. (\uline{Donc} \deleted{eff. admissible} alors $F$ est de poids $\in P'$, pour tt $P'\supset P$). Si \deleted{\unsure{$P\subset E$ et $P=\varnothing$, cela implique}} $P$ ne contient pas d'entier $\ge0$, cela implique $F=0$

% slanted
\margin{Attention, il n'est pas tautologique que si $F$ est eff. adm à poids $\in P$ et à poids $\in P'$, il soit : \deleted{\unsure{$P\cap P'$}} poids $\in P\cap P'$, \unsure{même si} $P\cap P'=\varnothing$}
% archive sheet number: 110
